% ═══════════════════════════════════════════════════════════════
% SPANISCH MUSTERLÖSUNG TEMPLATE
% ═══════════════════════════════════════════════════════════════
% Fachfarbe: #c41e3a (Karminrot)
% Lösungsfarbe: ROT (deutlich sichtbar)
%
% WICHTIG: Identisches Layout wie AB, nur mit roten Lösungen
% ═══════════════════════════════════════════════════════════════

\documentclass[11pt, a4paper]{article}

% ═══════════════════════════════════════════════════════════════
% SW-MODUS TOGGLE
% ═══════════════════════════════════════════════════════════════
% Setze \swmodustrue VOR \begin{document} für Schwarz-Weiß-Druck
\newif\ifswmodus
\swmodusfalse    % Standard: Farbe (auf true setzen für SW)

% ═══════════════════════════════════════════════════════════════
% PAKETE
% ═══════════════════════════════════════════════════════════════

\usepackage[utf8]{inputenc}
\usepackage[T1]{fontenc}
\usepackage[ngerman]{babel}
\usepackage{helvet}
\renewcommand{\familydefault}{\sfdefault}

\usepackage[margin=2cm]{geometry}
\usepackage{enumitem}
\usepackage{tabularx}
\usepackage{booktabs}
\usepackage{multirow}

\usepackage{xcolor}
\usepackage{framed}
\usepackage{tcolorbox}
\tcbuselibrary{skins}

\usepackage{fancyhdr}
\usepackage{lastpage}
\usepackage{needspace}          % Für \needspace (Seitenumbruch-Kontrolle)

% ═══════════════════════════════════════════════════════════════
% FARBEN
% ═══════════════════════════════════════════════════════════════

% Basis-Farben
\definecolor{spanisch-color}{HTML}{c41e3a}
\definecolor{grau-color}{HTML}{888888}
\definecolor{hellgrau-color}{HTML}{f5f5f5}
\definecolor{orange-color}{HTML}{b35c00}
\definecolor{loesung-color}{HTML}{c62828}

% SW-Farben (druckerfreundlich: keine Füllflächen, nur Linien)
\definecolor{sw-dark}{HTML}{000000}
\definecolor{sw-gray}{HTML}{666666}
\definecolor{sw-white}{HTML}{ffffff}

% Bedingte Farbzuweisung
\ifswmodus
    \colorlet{spanisch}{sw-dark}       % Rahmen/Text: schwarz
    \colorlet{grau}{sw-gray}           % Sekundärtext: dunkelgrau
    \colorlet{hellgrau}{sw-white}      % Hintergrund: weiß (keine Füllung)
    \colorlet{orange}{sw-dark}         % Hinweis-Rahmen: schwarz
    \colorlet{loesung}{sw-dark}        % Lösungen: schwarz
\else
    \colorlet{spanisch}{spanisch-color}
    \colorlet{grau}{grau-color}
    \colorlet{hellgrau}{hellgrau-color}
    \colorlet{orange}{orange-color}
    \colorlet{loesung}{loesung-color}
\fi

\newcommand{\fachfarbe}{spanisch}

% ═══════════════════════════════════════════════════════════════
% VARIABLEN - HIER AUSFÜLLEN
% ═══════════════════════════════════════════════════════════════

\newcommand{\thema}{[THEMA]}
\newcommand{\sequenz}{[SEQ]}
\newcommand{\block}{[BLOCK]}
\newcommand{\fach}{Spanisch}
\newcommand{\klasse}{7}
\newcommand{\maxpunkte}{[XX]}
\newcommand{\datum}{\today}

% ═══════════════════════════════════════════════════════════════
% KOPF- UND FUSSZEILE
% ═══════════════════════════════════════════════════════════════

\pagestyle{fancy}
\fancyhf{}
\renewcommand{\headrulewidth}{0.4pt}
\renewcommand{\footrulewidth}{0.4pt}

\lhead{\fach{} -- Klasse \klasse}
\chead{\textcolor{loesung}{\textbf{ML-\sequenz\block{} (MUSTERLÖSUNG)}}}
\rhead{\datum}

\lfoot{}
\cfoot{}
\rfoot{Seite \thepage\ von \pageref{LastPage}}

% ═══════════════════════════════════════════════════════════════
% BEFEHLE
% ═══════════════════════════════════════════════════════════════

% Punktebox
\newcommand{\punktebox}[1]{\hfill\fbox{\small \_/#1}}

% Lücke mit Lösung (ROT)
\newcommand{\loesung}[1]{\textcolor{loesung}{\textbf{#1}}}

% Leere Lücke (für Layout-Konsistenz)
\newcommand{\luecke}[1]{\rule{#1}{0.4pt}}

% Merkkasten (SW-kompatibel: schwebender Titel auf weißem Grund)
\newtcolorbox{merkkasten}[1][]{
    colback=hellgrau,
    colframe=\fachfarbe,
    colbacktitle=white,
    coltitle=\fachfarbe,
    boxrule=1pt,
    arc=0pt,
    left=8pt, right=8pt, top=8pt, bottom=8pt,
    fonttitle=\bfseries,
    attach boxed title to top left={yshift=-2mm, xshift=5mm},
    boxed title style={boxrule=0pt, colback=white},
    title=#1
}

% Vokabelkasten (kein Titel, daher unverändert)
\newtcolorbox{vokabelkasten}{
    colback=white,
    colframe=\fachfarbe,
    boxrule=0.5pt,
    arc=0pt,
    left=5pt, right=5pt, top=5pt, bottom=5pt
}

% Aufgabenumgebung (kein Seitenumbruch innerhalb)
\newenvironment{aufgabe}[1]{%
    \needspace{5\baselineskip}%
    \nopagebreak[4]%
    \vspace{10pt}
    \noindent\textbf{\textcolor{\fachfarbe}{Aufgabe #1}}
    \nopagebreak[4]%
    \vspace{5pt}
    \nopagebreak[4]%
    \begin{list}{}{\leftmargin=0pt}
    \item[]
}{%
    \end{list}
}

% Spanisch-Hervorhebung
\newcommand{\es}[1]{\textcolor{\fachfarbe}{\textbf{#1}}}

% ═══════════════════════════════════════════════════════════════
% DOKUMENT
% ═══════════════════════════════════════════════════════════════

\begin{document}

% ═══════════════════════════════════════════════════════════════
% KOPFBEREICH
% ═══════════════════════════════════════════════════════════════

\begin{center}
    {\Large\bfseries\textcolor{\fachfarbe}{Arbeitsblatt: \thema}}\\[0.3em]
    {\large\textcolor{loesung}{\textbf{--- MUSTERLÖSUNG ---}}}
\end{center}

\vspace{5pt}

\noindent
\begin{tabularx}{\textwidth}{|X|X|X|}
\hline
\textbf{Name:} \textcolor{loesung}{Musterschüler/in} &
\textbf{Klasse:} \klasse &
\textbf{Datum:} \datum \\
\hline
\end{tabularx}

\vspace{15pt}

% ═══════════════════════════════════════════════════════════════
% MERKKASTEN 1 (mit Lösungen)
% ═══════════════════════════════════════════════════════════════

\begin{merkkasten}[Merkkasten: [TITEL]]
\textbf{[THEMA]:}

\vspace{0.5em}

\begin{tabular}{ll}
\es{[SPANISCH 1]} & = \loesung{[DEUTSCH 1]} \\[0.8em]
\es{[SPANISCH 2]} & = \loesung{[DEUTSCH 2]} \\[0.8em]
\es{[SPANISCH 3]} & = \loesung{[DEUTSCH 3]} \\[0.8em]
\es{[SPANISCH 4]} & = \loesung{[DEUTSCH 4]} \\
\end{tabular}

\end{merkkasten}

\vspace{10pt}

% ═══════════════════════════════════════════════════════════════
% AUFGABE 1: Zuordnung (mit Lösungen)
% ═══════════════════════════════════════════════════════════════

\begin{aufgabe}{1}
Verbinde Spanisch und Deutsch. \punktebox{4}
\end{aufgabe}

\begin{center}
\begin{tabular}{rcl}
\es{[ES\_1]} & \loesung{---} & [DE\_A] \loesung{(= [RICHTIGE ZUORDNUNG])} \\[0.8em]
\es{[ES\_2]} & \loesung{---} & [DE\_B] \loesung{(= [RICHTIGE ZUORDNUNG])} \\[0.8em]
\es{[ES\_3]} & \loesung{---} & [DE\_C] \loesung{(= [RICHTIGE ZUORDNUNG])} \\[0.8em]
\es{[ES\_4]} & \loesung{---} & [DE\_D] \loesung{(= [RICHTIGE ZUORDNUNG])} \\
\end{tabular}
\end{center}

\vspace{15pt}

% ═══════════════════════════════════════════════════════════════
% AUFGABE 2: Lückentext (mit Lösungen)
% ═══════════════════════════════════════════════════════════════

\begin{aufgabe}{2}
Ergänze die Lücken. \punktebox{4}
\end{aufgabe}

\begin{vokabelkasten}
\textit{Wortliste:} [WORT\_1] | [WORT\_2] | [WORT\_3] | [WORT\_4]
\end{vokabelkasten}

\vspace{0.5em}

\noindent
a) \es{¡\loesung{Buenos días}!} -- Das sagt man morgens. \\[0.8em]
b) \es{¡\loesung{Buenas noches}!} -- Das sagt man abends. \\[0.8em]
c) \es{¡\loesung{Adiós}!} -- Das sagt man zum Abschied. \\[0.8em]
d) \es{¿\loesung{Cómo te llamas}?} -- Das fragt man nach dem Namen. \\

\vspace{15pt}

% ═══════════════════════════════════════════════════════════════
% AUFGABE 3: Dialog vervollständigen (mit Lösungen)
% ═══════════════════════════════════════════════════════════════

\begin{aufgabe}{3}
Vervollständige den Dialog. \punktebox{6}
\end{aufgabe}

\noindent
\textbf{A:} \es{¡Hola! ¿Cómo te llamas?}

\textbf{B:} \loesung{Me llamo [Name]. / Soy [Name].}

\textbf{A:} \es{¿De dónde eres?}

\textbf{B:} \loesung{Soy de Alemania. / Soy de [Ort].}

\textbf{A:} \es{¡Mucho gusto!}

\textbf{B:} \loesung{¡Igualmente! / ¡Mucho gusto!}

\vspace{15pt}

% ═══════════════════════════════════════════════════════════════
% AUFGABE 4: Freie Produktion (Beispiellösung)
% ═══════════════════════════════════════════════════════════════

\begin{aufgabe}{4}
Schreibe einen kurzen Dialog (mind. 4 Zeilen). \punktebox{6}
\end{aufgabe}

\textbf{Beispiellösung:}

\loesung{A: ¡Hola! ¿Qué tal?}

\loesung{B: Bien, gracias. ¿Y tú?}

\loesung{A: Muy bien. ¿Cómo te llamas?}

\loesung{B: Me llamo María. ¡Hasta luego!}

\loesung{A: ¡Adiós!}

\vspace{0.5em}
\textit{\textcolor{grau}{Bewertung: Begrüßung (1P), Frage (1P), Antwort (2P), Verabschiedung (1P), sprachliche Korrektheit (1P)}}

% ═══════════════════════════════════════════════════════════════
% PUNKTEÜBERSICHT
% ═══════════════════════════════════════════════════════════════

\vspace{20pt}
\hrule
\vspace{10pt}

\begin{flushright}
\textbf{Punkte:} \loesung{\maxpunkte} / \maxpunkte
\end{flushright}

\end{document}
